% explainer-source-hash: 378da11993b49d26f51eace8aae858845ac55c6ff180b36302677526a7ffcf01
% explainer-source-commit: b46124dd2204fce0b088f79ad7cd690afbbcf024
% explainer-generated: 2026-07-16
% Regenerate with /explain-all (see WIRING.md). Do not edit this stamp by hand:
% it is how the toolchain knows whether this document still matches the code.
\documentclass[11pt,a4paper]{article}
\input{preamble}
\begin{document}

\begin{center}
{\LARGE\bfseries\color{inkblue} Morse Code Converter\par}
\vspace{4pt}
{\large\color{accent} The Brief\par}
\vspace{6pt}
{\small\itshape A high-level but complete account of what the project is, how it is
built, and where it is weak. Terms are defined in the glossary on the last page.}
\end{center}
\vspace{6pt}
\hrule
\vspace{10pt}

% =====================================================================
\section{The project in one paragraph}

A desktop application in Python that converts text to International Morse Code and
back, live, as you type, in a two-tab window. It can play the current Morse code
aloud as real 700~Hz beeps while flashing each dot and dash on screen in exact time
with the sound. It is three Python files, roughly 600 lines including comments, and
it has \textbf{no third-party dependencies}: the interface, the sound generation and
the playback all run on Python's standard library alone. There is no database, no
network access, and no configuration.

It began as a practice exercise in dictionary lookups and grew into an application.
That history is visible in the code, and mostly in a good way.

% =====================================================================
\section{Architecture}

The organising decision is that \textbf{only one of the three files knows a screen
exists}. The conversion logic and the audio logic can each be imported and used with
no window ever opening, which is what makes the logic testable and let the audio be
rewritten without touching the converter.

\begin{figure}[H]
\centering
\begin{tikzpicture}[node distance=10mm]
  \node[box] (gui) {\file{gui.py}\\[2pt] \scriptsize tkinter window: two tabs, live text\\ \scriptsize boxes, playback canvas, buttons};
  \node[box, below left=16mm and 4mm of gui] (conv) {\file{morse\_converter.py}\\[2pt] \scriptsize 57-entry alphabet table\\ \scriptsize \code{text\_to\_morse} / \code{morse\_to\_text}\\ \scriptsize \textit{imports nothing at all}};
  \node[box, below right=16mm and 4mm of gui] (audio) {\file{morse\_audio.py}\\[2pt] \scriptsize standard Morse timing\\ \scriptsize sine-wave WAV generation\\ \scriptsize runtime backend detection};
  \node[extbox, below=12mm of audio] (os) {system audio player\\[2pt] \scriptsize \code{aplay}, \code{paplay},\\ \scriptsize \code{afplay}, \code{ffplay},\\ \scriptsize or \code{winsound}};
  \draw[flow] (gui) -- node[lbl,left] {imports} (conv);
  \draw[flow] (gui) -- node[lbl,right] {imports} (audio);
  \draw[dflow] (audio) -- node[lbl,right] {spawns} (os);
\end{tikzpicture}
\caption{Module dependencies. Nothing points upward, and the two lower modules do
not import each other in any real code path.}
\end{figure}

\subsection{\file{morse\_converter.py}: the alphabet}

A hardcoded dictionary maps 57 characters (26 letters, 10 digits, the space, and 20
punctuation marks) to their Morse patterns. A one-line dictionary comprehension
inverts that same table to build the decoder, so there is exactly one place to edit a
mapping. An excerpt, exactly as it appears in the source (the Deep Dive prints all 57
entries):

\begin{lstlisting}[language=Python]
MORSE_CODE = {
    "A": ".-", "B": "-...", "C": "-.-.", "D": "-..", "E": ".",
    # ... through N ...
    "O": "---", "P": ".--.", "Q": "--.-", "R": ".-.", "S": "...",
    # ... through Z, then the ten digits ...
    "0": "-----", "1": ".----", "2": "..---", "3": "...--",
    " ": "/",
    "!": "-.-.--", "?": "..--..", ".": ".-.-.-", ",": "--..--",
    # ... twenty punctuation marks in total, including "-": "-....-" ...
}
\end{lstlisting}

So \code{text\_to\_morse("SOS")} looks up S, O and S in turn and returns
\code{... --- ...}, the canonical Morse message: three dots, three dashes, three
dots.

\code{text\_to\_morse} uppercases the input (this is the entire ``case-insensitive''
feature) and looks up each character, raising a \code{ValueError} that names the
offending character rather than letting a low-level \code{KeyError} escape.
\code{morse\_to\_text} reverses the process. Running the file directly still offers
the original one-line command-line prompt, though it can only encode, not decode.

\begin{keyidea}{The row that does the quiet work}
The space character is treated as an ordinary table entry whose Morse ``pattern'' is
a forward slash (\code{" ": "/"}). Because \code{text\_to\_morse} joins per-character
codes with single spaces, every word boundary in the output is literally the three
characters \code{"\ /\ "}. That regularity is what makes decoding, and the playback
timing, straightforward to parse.
\end{keyidea}

\subsection{\file{morse\_audio.py}: real sound from nothing}

Morse code is a pattern of \emph{durations}, not just of symbols, and the standard
fixes them all as multiples of one base unit (80~ms here):

\begin{table}[H]
\centering
\begin{tabular}{lll}
\toprule
dot & 1 unit of tone & gap between letters: 3 units of silence \\
dash & 3 units of tone & gap between words: 7 units of silence \\
gap inside a letter & 1 unit of silence & \\
\bottomrule
\end{tabular}
\caption{The standard ratios, each a named constant in the source. Silence carries
information: without three distinct gap lengths a listener cannot tell where one
letter ends and the next begins.}
\end{table}

\code{build\_timing\_sequence} turns a Morse string into an explicit ordered list of
\code{(kind, units)} events, emitting gaps only \emph{between} elements so a message
never ends in trailing silence.

Sound itself is generated from first principles: \code{math.sin} produces a sine
wave, \code{struct} packs the samples into raw bytes, and the standard library's
\code{wave} module writes them to a WAV file. This is the whole reason the
dependency list is empty. The project's own README records that it originally
shipped visual-only playback because there ``seemed to be no dependency-free way to
make a sound'', and that this was simply wrong.

Playing the file is the harder half, because Python has no cross-platform way to
reach a speaker. \code{detect\_playback\_backend} therefore \textbf{asks the machine
rather than assuming}: it tries \code{winsound} on Windows, then checks for
\code{aplay}, \code{paplay}, \code{afplay} and \code{ffplay} with \code{shutil.which}
in that order, and treats ``nothing found'' as a legitimate answer. On this machine
it returns \code{aplay}. If it returns nothing, the app says so plainly in the
playback panel and keeps the visual flash working alone.

Tones are launched with \code{subprocess.Popen}, which starts the player and returns
immediately rather than waiting for it to finish. This is load-bearing: if playback
blocked for the 240~ms of a dash, the window would freeze and the picture would fall
behind the sound.

\subsection{\file{gui.py}: the window}

A hand-styled dark tkinter interface. One factory function builds both tabs, handed a
different conversion function each time, so identical machinery serves both
directions. Each input box re-converts its entire contents on every \code{<KeyRelease>}
event, which is why there is no submit button; conversion therefore runs constantly on
invalid half-typed input, and a caught \code{ValueError} becomes a small inline
message instead of a crash. Audio is detected and the two tone files generated at
startup, before anything is drawn, so no delay is audible during playback.

% =====================================================================
\section{The idea worth taking away}

\begin{figure}[H]
\centering
\begin{tikzpicture}[node distance=9mm]
  \node[box] (morse) {the current Morse string};
  \node[box, below=of morse] (seq) {\code{build\_timing\_sequence()}\\[2pt] \scriptsize called \textbf{once}};
  \node[box, below=of seq] (list) {one list of \code{(kind, units)} events};
  \node[box, below left=11mm and -4mm of list] (eye) {canvas flash\\[2pt] \scriptsize driven by \code{after()}};
  \node[box, below right=11mm and -4mm of list] (ear) {tone playback\\[2pt] \scriptsize \code{Popen}, non-blocking};
  \draw[flow] (morse) -- (seq);
  \draw[flow] (seq) -- (list);
  \draw[flow] (list) -- node[lbl,left] {the eyes} (eye);
  \draw[flow] (list) -- node[lbl,right] {the ears} (ear);
\end{tikzpicture}
\caption{One timeline, two consumers.}
\end{figure}

\begin{keyidea}{Why this matters}
The visual flash and the audio read from the \textbf{same} event list, so they cannot
drift apart. The obvious alternative, hand-tuning a visual delay and an audio delay
separately, is exactly what this project used to do, and the README records the
result: a dot-to-dash ratio of roughly 1:1.9 instead of the correct 1:3, and one flat
approximate gap where the standard requires three different ones. Deriving both
outputs from a single source of truth is the fix, and it is the most transferable
lesson in the project.
\end{keyidea}

Two other decisions are worth naming. \textbf{Detecting the audio backend at runtime}
rather than inferring it from the operating system means a machine with no player
installed gets an honest message instead of a crash or a silent no-op. And
\textbf{keeping the conversion functions free of any tkinter import} is what allowed
the entire audio subsystem to be rewritten, twice, without the converter changing.

% =====================================================================
\section{What I verified, and what I could not}

Every behavioural claim below was checked by running the code, not read from the
README:

\begin{itemize}
\item The 57-entry table inverts \textbf{losslessly}: no two characters share a Morse
      pattern, so the decoder loses nothing. Notably, nothing in the code
      \emph{checks} this; it is correct by the table's contents, not by any guard.
\item The round-trip property \code{morse\_to\_text(text\_to\_morse(s)) == s.upper()}
      held for every input tested: mixed case, digits, punctuation, a literal
      \code{/}, leading and trailing spaces, doubled spaces, all-whitespace.
\item \textbf{Timing is exact.} Rendering \code{"SOS"} gave 27 units, and the
      resulting WAV measured 95{,}256 frames at 44{,}100~Hz, precisely 2160~ms,
      matching $27 \times 80$~ms computed by hand from the standard ratios.
      \code{"SOS SOS"} gave 61 units and 215{,}208 frames, exactly 4880~ms, so the
      7-unit word gap is accounted for to the sample.
\item \code{detect\_playback\_backend()} returns \code{'aplay'} here, as documented.
\end{itemize}

I could not run the GUI (no display in this session), so statements about on-screen
behaviour are read from the code. The \code{winsound}, \code{afplay} and
\code{paplay} paths need Windows, macOS and PulseAudio respectively and were not
exercised; the README states this limitation itself.

% =====================================================================
\section{Honest assessment}

The good is real. The module separation is genuine rather than nominal, the single
event list is the right answer to a problem the project first got wrong, the runtime
backend detection with a working fallback is more care than a practice project
usually receives, and generating correct audio from \code{math.sin} and \code{wave}
is a well-earned way to keep the dependency list empty. The README is unusually
candid about the project's own missteps.

The weaknesses, in order of how much they matter:

\begin{honest}{No automated tests}
The project's correctness lives almost entirely in a 57-entry table, and there is no
test file. Every claim in the section above had to be re-established by hand. A
silently wrong table entry is exactly the defect that manual spot-checks miss and a
test would catch instantly. The README names this first among what it would do
differently, which is the right instinct.
\end{honest}

\begin{honest}{Play does not respect the tab you are looking at}
\code{\_current\_tab\_state()} is misnamed twice over: it returns a string, not a tab
state, and it never asks which tab is selected. It reads the Text-to-Morse tab's
output first and only falls back to the other tab. So typing in the first tab,
switching to the second, typing there and pressing Play replays the \emph{first}
tab's content. The README acknowledges the fixed priority honestly, but the behaviour
is still surprising.
\end{honest}

\begin{honest}{Two contracts for the same malformed input}
\code{morse\_to\_text} raises a \code{ValueError} naming a bad symbol, while
\code{build\_timing\_sequence} silently discards anything that is not a dot or dash.
I confirmed that \code{build\_timing\_sequence(".- XYZ -...")} drops \code{XYZ}
without complaint yet still emits the gaps around it. The GUI masks this, but the two
halves of the library genuinely disagree.
\end{honest}

\begin{honest}{Redundant decoding logic, with a rationale that does not hold}
\code{morse\_to\_text} handles word gaps twice: it splits the input on \code{" / "},
and \code{"/"} is \emph{also} an ordinary decodable table entry meaning space. A
naive one-line decoder that never splits on \code{" / "} produced identical output on
every input I tested. The README defends the split by arguing that ``splitting
naively on \code{/} alone would have left stray spaces'', but the alternative it
overlooks is not splitting on the slash at all, and simply decoding it like any other
symbol. Not a bug; redundant complexity with a stale justification.
\end{honest}

Smaller items: \code{\_safe\_text\_to\_morse} and \code{\_safe\_morse\_to\_text} are
pure pass-through wrappers whose \code{\_safe\_} prefix promises a safety they do not
provide (the real error handling is elsewhere). The multi-line input boxes reject
newlines, since \code{"\textbackslash n"} is not in the table, so pressing Enter
blanks the output with an inline error; it is handled gracefully but the widget
invites input the alphabet cannot represent. \code{render\_morse\_to\_wav} is unused
by the application and is really a development instrument. And because tone processes
are never collected, each finished player leaves a short-lived zombie process until
the app exits.

None of these is a crash, a security issue, or a wrong answer. They are the ordinary
residue of a script that grew into an application, and the distance between this
code and its documentation is small.

% =====================================================================
\section*{Glossary}
\addcontentsline{toc}{section}{Glossary}

\begin{description}[leftmargin=0pt, style=nextline, itemsep=3pt]
\item[Morse code] Representing letters and digits as sequences of short signals
  (dots) and long signals (dashes), designed for telegraph wires where the only
  controllable thing was whether current flowed. ``International'' (ITU) Morse is the
  modern standardised form.
\item[Dictionary (\code{dict})] A container mapping keys to values, with instant
  lookup by key regardless of size, because it uses a hash table internally rather
  than scanning entries.
\item[Module] A single Python file. One module can \emph{import} another to use what
  is defined inside it.
\item[GUI / tkinter / widget] A graphical user interface is the visible, clickable
  part of a program. \emph{tkinter} is Python's built-in GUI toolkit; a \emph{widget}
  is any element within it (window, button, text box, canvas).
\item[Event loop] The endless loop at a GUI's heart that waits for something to
  happen (a key, a click, a timer) and calls the function registered for it.
  \code{root.after(ms, fn)} asks it to call \code{fn} later without blocking.
\item[Callback] A function handed to something else to be called later, when an event
  occurs.
\item[Exception / raising / \code{ValueError}] When a function cannot do its job it
  \emph{raises} an exception, stopping immediately and handing an error up to its
  caller, who may \emph{catch} it. \code{ValueError} conventionally means ``this
  value is unusable''.
\item[PCM / sample rate] Digital audio stores sound by measuring air pressure
  thousands of times per second. Each measurement is a \emph{sample}; the rate here
  is 44{,}100 per second, the audio-CD rate. The raw list of numbers is \emph{PCM}.
\item[Sine wave / frequency / hertz] A pure tone is a value swinging smoothly up and
  down at a fixed rate. That rate, in hertz, is the frequency, perceived as pitch.
  This project uses 700~Hz.
\item[WAV / the \code{wave} module] A WAV file is raw PCM samples with a small header
  declaring sample rate, channel count and sample size. Python's standard library
  writes it, which is why no audio dependency is needed.
\item[Process / \code{subprocess.Popen} / blocking] A process is a running program.
  \code{Popen} launches one and returns immediately; \code{run} and \code{wait}
  instead \emph{block}, freezing the caller until it finishes.
\item[\code{shutil.which}] Answers ``would this command actually run on this machine,
  and where does it live?'' by searching \code{PATH}. Returns \code{None} if not.
\item[Standard library] The code that ships with Python itself, needing no
  installation. Everything this project uses comes from it.
\end{description}

\end{document}
